\section{Contexte : ce que le projet savait d\'ej\`a}

\subsection{Batching et virtual loss}

La recherche regroupe d\'ej\`a plusieurs feuilles dans une seule inf\'erence.
Chaque descente pose un \emph{virtual loss} de type Leela Chess Zero : un
compteur \code{n\_in\_flight} incr\'ement\'e le long du chemin, qui n'entre que
dans le d\'enominateur du terme d'exploration et laisse $Q$ intact. Le noyau
mono-worker collecte jusqu'\`a \code{batch\_size} feuilles, s'arr\^ete au
premier n\oe{}ud d\'ej\`a r\'eserv\'e, puis appelle l'\'evaluateur une fois. Le
balayage de batching du 14 septembre 2026 avait fix\'e le lot de production
\`a 8, meilleur compromis entre d\'ebit et qualit\'e.

\subsection{Vagues multic\oe{}ur}

Le chantier multic\oe{}ur a remplac\'e la collecte mono-c\oe{}ur par des
collecteurs persistants travaillant sur un arbre partag\'e, avec r\'eservation
atomique des feuilles, table de transposition par snapshots, garde RAII et
budget exact. Le rapport de r\'esultats \code{2026-09-16-multicore-waves-results.md}
mesure :

\begin{itemize}[leftmargin=1.2em,itemsep=2pt]
  \item des gains de $+6$ \`a $+22\,\%$ de d\'ebit selon la position entre 1 et
  8 workers, avec un plafond d\`es deux workers ;
  \item un p95 de latence sous $+5\,\%$ partout, et m\^eme am\'elior\'e en
  finale ;
  \item une qualit\'e non inf\'erieure sur 2\,500 puzzles (delta $-0,28$ point,
  IC95 $[-0,8\,;+0,24]$) ;
  \item l'activation de \code{MCTS\_WORKER\_COUNT = 8} dans \code{uci.py}.
\end{itemize}

\subsection{R\'epartition par phase}

Les chronom\^etrages internes (\code{phases.json}, 700 simulations, lot 8)
montrent une domination \'ecrasante de l'\'evaluateur :

\begin{table}[h]
\centering
\small
\caption{Part de l'\'evaluateur dans le temps mur, 700 simulations.}
\label{tab:phases}
\begin{tabular}{lrrrr}
\toprule
Position & Total (ms) & \'Evaluateur (ms) & Part & CPU + barri\`ere \\
\midrule
Ouverture, w1 & 923,8 & 903,2 & 97,8\,\% & 20,6\,ms \\
Milieu, w1 & 2\,063,5 & 2\,038,1 & 98,8\,\% & 25,4\,ms \\
Finale, w1 & 1\,447,5 & 1\,426,4 & 98,5\,\% & 21,1\,ms \\
Finale, w8 & 1\,300,4 & 1\,264,0 & 97,2\,\% & 36,4\,ms \\
Milieu, w8 & 2\,335,3 & 2\,280,8 & 97,7\,\% & 54,5\,ms \\
Ouverture, w8 & 853,3 & 824,2 & 96,6\,\% & 29,1\,ms \\
\bottomrule
\end{tabular}
\end{table}

Le travail CPU de recherche, coups compris, vaut 29 \`a 36 microsecondes par
simulation, soit 1 \`a 1,5\,\% du temps. La barri\`ere des vagues co\^ute 1,5 \`a
2,3\,\%. Autrement dit, tout gain durable ne peut venir que de l'appel
r\'eseau.

\subsection{Le mod\`ele \texorpdfstring{$F + M \cdot N$}{F + M x N}}

Le balayage de lots du 14 septembre 2026, mesur\'e dans une seule session,
montre que le d\'ebit plafonne d\`es le lot 8 : demander 64 positions n'en fait
\'evaluer que 5 \`a 10, parce que l'arbre ne fournit pas plus de feuilles
utiles. En ajustant $T(N) = F + M \cdot N$ sur le lot 1 et le lot 8, on trouve
un co\^ut fixe $F$ d'environ 1 \`a 2\,ms et un co\^ut par position $M$ de 0,9 \`a
1,9\,ms selon la position, d'o\`u un plafond du batching seul de
$\times 1{,}6$ \`a $\times 3{,}4$.

Ce mod\`ele d\'ecrivait bien les gains observ\'es, mais il laissait une
incoh\'erence : le self-play, avec un lot plus grand mais surtout \emph{fixe},
obtenait 0,38\,ms par position, incompr\'ehensible dans un mod\`ele o\`u $M$
devrait \^etre une propri\'et\'e du r\'eseau.

\subsection{L'\'etude de co\^ut et ses hypoth\`eses}

L'\'etude \code{2026-09-19-cout-calcul-cpu-gpu.md} avait agr\'eg\'e les mesures
et compar\'e le moteur \`a Leela Chess Zero. Elle estimait l'\'ecart \`a un
facteur 15 \`a 200 selon quatre axes : taille de lot ($\times 3$ \`a $8$),
ex\'ecution asynchrone ($\times 1{,}2$ \`a $1{,}5$), backend r\'eseau
($\times 2$ \`a $4$) et cache (\`a $4$). Elle notait surtout que le premier
facteur, la plomberie de l'appel, \'etait d\'ej\`a prouv\'e par le banc
self-play, et que le lot d'un seul arbre plafonnant vers 7, le gain total
r\'ealiste \'etait de $\times 3$ \`a $\times 6$, pas $\times 30$.

C'est cette piste, et non le r\'eseau, que le plan
\code{2026-09-19-cout-calcul-optimisations.md} a mise en premier, avec un banc
hors arbre comme \'etape pr\'ealable.

\subsection{Le r\'egime r\'eel du bot}

Le bot n'appelle pas \code{step\_analysis} une fois pour 700 simulations : il
d\'ecoupe en tranches de \code{BATCH\_SIZE = 64} pour garder la main sur
l'horloge, avec arbre et pool persistants. Le rapport multic\^ur mesurait d\'ej\`a
que le gain des workers d\'ependait de la tranche : $+5$ \`a $+27\,\%$ \`a 64
simulations, mais $-19\,\%$ \`a 8, o\`u les arbres naissants provoquent environ
970 collisions pour 800 simulations. Toutes les campagnes de ce rapport
mesurent donc aussi aux tranches r\'eelles, en plus des appels longs.
